\section{Introdução}
\label{sec:intro}

A proliferação de dispositivos de \emph{Internet of Things} (IoT) e
\emph{Industrial IoT} (IIoT) ampliou drasticamente a superfície de
ataque das redes corporativas e residenciais. Sistemas de detecção de
intrusão (\emph{Intrusion Detection Systems}, IDS) baseados em
aprendizado de máquina tornaram-se a abordagem dominante nesse cenário,
em particular quando operam sobre a telemetria coletada diretamente nos
sensores (temperatura, posição, sinal de porta, leituras Modbus). Essa
telemetria é o sinal mais próximo do fenômeno físico monitorado, o que
a torna um ponto natural de defesa, mas também o mais restrito em
recursos de processamento.

Detectar intrusões nessa camada impõe dois requisitos que competem entre
si no \emph{deployment} em \emph{edge}. O primeiro é eficácia preditiva
alta o suficiente para que falsos positivos não inundem o operador; o
segundo é custo computacional baixo o bastante para caber em hardware
embarcado de recursos limitados, longe da memória e da energia de um
servidor. Modelos baseados em árvores ou em redes profundas atingem o
primeiro requisito, mas a um custo de memória e inferência que tensiona
o segundo, e essa tensão entre eficácia e custo permanece sem uma
solução de baixo custo bem caracterizada nesse domínio.

A literatura de IDS sobre a telemetria do TON\_IoT não resolve essa
tensão porque trata a acurácia quase-perfeita como meta a otimizar, não
como sintoma a investigar. Estudos recentes reportam acurácia de até
$100\%$ em bases isoladas e macro-F1 acima de $0{,}98$ na telemetria sem
qualquer auditoria de vazamento de
dados~\cite{sharrab2025analysis,alotaibi2023stacking,benaddi2025near}.
As críticas que de fato diagnosticam vazamento atuam apenas sobre o
\emph{dataset} de rede, descartando colunas identificadoras como IP e
porta~\cite{shaikhanova2025audit,dharini2026efficient}, colunas que
inexistem na telemetria; a opção de baixo custo computacional do
compromisso e a auditoria da própria telemetria permanecem ambos sem
mapeamento.

As redes neurais sem peso (\emph{Weightless Neural Networks},
WNN)~\cite{aleksander1984wisard,aleksander2009brief} oferecem, em
princípio, um caminho para essa tensão de baixo custo não mapeada. Uma
WNN substitui os produtos peso-entrada das redes convencionais por
consultas $\mathcal{O}(1)$ em \emph{RAMs} indexadas por padrões binários
da entrada, de modo que o custo de inferência é limitado pelo acesso a
tabela e independe da ``profundidade'' do modelo. Essa propriedade é
exatamente a que poderia manter um IDS de \emph{edge} barato sem ceder
eficácia preditiva, trocando aritmética em ponto flutuante por leitura
de memória. É essa hipótese, e não um resultado assumido, que motiva
avaliar a família WiSARD de WNNs no problema de telemetria.

Este trabalho avalia, sob esse duplo critério, a \textbf{família WiSARD}
de redes neurais sem peso contra oito \emph{baselines} clássicos sobre o
\textbf{TON\_IoT}~\cite{alsaedi2020toniot}, \emph{benchmark} amplamente
adotado para IDS em IoT/IIoT cujo paper de referência, Alsaedi et al.\
(IEEE Access, 2020), estabelece a linha de base que reproduzimos.
Avaliados em sete bases de telemetria e em uma base consolidada sob o
protocolo do paper, WiSARD e ClusWiSARD lideram a família a
F1 \emph{weighted} $\approx 0{,}82$ \emph{per-device} (WiSARD $0{,}824$,
ClusWiSARD $0{,}822$), a cerca de um ponto percentual de Random Forest.
A família sem peso, portanto, ocupa a opção de baixo custo do
compromisso quase sem abrir mão de eficácia, e a auditoria que
conduzimos no caminho revela um \emph{leak} categórico determinístico
que ajuda a explicar a acurácia quase-perfeita aceita pela literatura.

\paragraph{Pergunta de pesquisa.} A família WiSARD, projetada
explicitamente para baixa demanda computacional, consegue se aproximar
da eficácia preditiva dos \emph{baselines} clássicos no TON\_IoT
quando avaliada \emph{sob o mesmo protocolo experimental do paper de
referência}?

\paragraph{Contribuições.} Este relatório consolida os seguintes
resultados:

\begin{enumerate}[leftmargin=*,itemsep=0pt]
  \item \textbf{Reprodução \emph{paper-faithful} dos oito
        \emph{baselines}} de Alsaedi 2020 (LR, LDA, kNN, RF, CART, NB,
        SVM e LSTM) sob o protocolo descrito em
        sua Seção VI--A (\emph{split} 80/20, $k=4$ \emph{StratifiedKFold}
        no \emph{train}, métrica \emph{F1 weighted} adotada após
        comparação empírica). Os números obtidos batem com o paper
        dentro de $|\Delta| < 0{,}15$ em 91\% dos pares
        (modelo, base).
  \item \textbf{Avaliação de cinco modelos da família WiSARD}
        (WiSARD~\cite{aleksander1984wisard},
        ClusWiSARD~\cite{cardoso2016cluswisard},
        BloomWiSARD~\cite{santiago2020bloomwisard},
        BTHOWeN~\cite{susskind2022bthowen},
        ULEEN~\cite{susskind2023uleen}) sob o \emph{mesmo} protocolo
        experimental aplicado aos \emph{baselines} do paper.
  \item \textbf{Auditoria de qualidade do dataset TON\_IoT}, em que
        identificamos três problemas estruturais nos CSVs públicos:
        \emph{leak} categórico determinístico em duas bases e divergência
        de tamanho contra a versão usada nas Tabelas 10--13 do paper
        (ambos não discutidos no paper), além do \emph{leak} temporal
        que o paper já descarta e que aqui caracterizamos como
        determinístico.
  \item \textbf{Extensões na biblioteca \texttt{wisardpkg}}. Estendemos a
        \texttt{wisardpkg} com cinco termômetros ajustados aos dados
        (\emph{Distributive}, \emph{Gaussian}, \emph{Exponential},
        \emph{Logarithmic} e a alocação \emph{Non-uniform}), avaliados
        como seis variantes junto ao termômetro \texttt{linear} clássico
        na análise de sensibilidade. Adicionamos ainda a métrica
        \texttt{deployedSizeBytes}, que normaliza o custo de memória das
        variantes sem peso em um eixo único.
\end{enumerate}

\paragraph{Estrutura do documento.}
A Seção~\ref{sec:related} apresenta o paper de referência e os modelos
da família WiSARD. A Seção~\ref{sec:dataset} reporta a auditoria de
qualidade do TON\_IoT: os \emph{leaks} categóricos e temporais
identificados e a divergência de tamanho entre os CSVs públicos e o
paper. A Seção~\ref{sec:metodologia} formaliza o protocolo
\emph{paper-faithful} usado para todas as comparações.
A Seção~\ref{sec:reproducao} reporta a reprodução dos oito
\emph{baselines}, validando o pipeline com sanity checks derivados das
Tabelas 10--13 do paper. A Seção~\ref{sec:wisard} apresenta o
\emph{sweep} de hiperparâmetros dos cinco modelos da família WiSARD.
A Seção~\ref{sec:comparacao} consolida as três visões
lado-a-lado: paper, nossa reprodução e WiSARD-family.
A Seção~\ref{sec:discussao} discute limitações honestas da reprodução,
e a Seção~\ref{sec:conclusao} encerra com as conclusões.
